\section{Related work}
\label{sec:related}

This section positions the present work relative to (i) supercompilation and program transformation, (ii) cyclic proof theory, (iii) sequent calculi for dependent types, and (iv) proof-theoretic treatments of induction.

\subsection{Supercompilation}

Supercompilation is a program transformation technique based on symbolic evaluation (driving), case analysis, generalization, and fold/refold steps \cite{turchin-1986-supercompiler,sorensen-gluck-1995-supercompilation}.
The key operations are:
\begin{itemize}
\item \textbf{Driving}: unfolding function definitions and performing symbolic reduction
\item \textbf{Splitting}: case analysis on unknown values (neutrals)
\item \textbf{Generalization}: abstracting common patterns to enable folding
\item \textbf{Folding}: recognizing repeated configurations and creating recursive calls
\end{itemize}

\paragraph{Termination control.}
A central challenge in supercompilation is ensuring termination. 
Classical approaches use \emph{whistles} based on homeomorphic embedding or well-quasi-orderings \cite{kruskal-1960-wqo,leuschel-2002-homeomorphic,sorensen-gluck-jones-1996-positive}.
These approaches analyze configuration histories and force generalization when a dangerous pattern is detected.
Our cyclic proof framework replaces the whistle with a \emph{global trace condition}: we allow arbitrary folding but require that every cycle passes through a progress event (a split on a structurally smaller inductive component).

\paragraph{Supercompilation and partial evaluation.}
Early work connects supercompilation to partial evaluation \cite{jonesbook-1996-partial-evaluation}.
Our contribution extends this connection to dependent types: we show that supercompilation \emph{is} proof search in a typed setting, with configurations as sequents and the memo table as a cyclic proof graph.

\paragraph{Prior work on dependent types.}
To our knowledge, this is the first formulation of supercompilation for a dependently typed core language.
Existing work on supercompilation focuses on first-order functional languages (Refal, Scheme, Haskell subsets) or simple typed settings \cite{sorensen-gluck-1995-supercompilation,jonesbook-1996-partial-evaluation}.
The dependent type context introduces new challenges:
\begin{itemize}
\item Types depend on terms, so driving must respect typing constraints
\item The context $\Gamma$ changes during proof search (extended in branches)
\item Observational equivalence is more subtle (CIU for dependent types)
\end{itemize}

Our sequent calculus formulation addresses these challenges by making the typing context explicit in every sequent and using type-directed driving rules.

\subsection{Cyclic proof theory}

Cyclic proof systems represent potentially infinite derivations using finite directed graphs with back-edges, validated by a global soundness condition \cite{brotherston-2006-cyclic,brotherston-simpson-2011-sequent}.

\paragraph{Concrete cyclic systems.}
Brotherston and Simpson develop cyclic proofs for first-order logic with inductive definitions, using a global trace condition to ensure soundness.
Their work establishes that cyclic proofs are as expressive as standard proofs with explicit induction rules, but offer computational advantages (the induction scheme emerges from the cycle structure).

\paragraph{Abstract cyclic frameworks.}
Afshari and Wehr develop a modern, highly abstract account of cyclic proof checking based on trace conditions, activation algebras, and categorical structure \cite{afshari-wehr-2024-abstract-cyclic}.
Our work instantiates their abstract framework in a concrete setting: CoC with inductives, CBN operational semantics, and observational equivalence.

\paragraph{Cyclic proofs as normal forms.}
Our contribution extends cyclic proof theory in a new direction: we use cyclic proofs not just as proof objects, but as \emph{normal forms for program transformations}.
The correspondence with supercompilation shows that cyclic proofs are stable under driving, splitting, and folding---operations that would be difficult to justify in a traditional natural deduction setting.

\subsection{Sequent calculi for dependent types}

\paragraph{Bidirectional type checking.}
Bidirectional type systems \cite{pierce-turner-2000-local-type-inference} separate synthesis (inferring types) from checking (validating against known types).
This discipline is standard in dependently typed proof assistants (Agda, Rocq, Lean).
Our sequent calculus uses a similar decomposition, but with a crucial difference: we use \emph{focused} sequents that separate asynchronous (deterministic) from synchronous (choice-bearing) phases.

\paragraph{Focused sequent calculi.}
Focusing \cite{andreoli-1992-logic-programming,miller-2004-forum} is a proof search discipline that reduces nondeterminism by eagerly applying invertible rules (asynchronous phase) and carefully controlling choice points (synchronous phase).
Focusing has been applied to linear logic, classical logic, and intuitionistic logic, but \textbf{to our knowledge this is the first focused sequent calculus for CoC with inductives}.

The key novelty is the connection to supercompilation: the asynchronous phase corresponds to driving (deterministic unfolding), and the synchronous phase corresponds to splitting and folding (the essential choice points).

\paragraph{Sequent calculi for CIC.}
Standard presentations of CIC use natural deduction \cite{coquand-huet-1988-coc,paulin-mohring-1993-inductive}.
Herbelin develops a sequent calculus for the Calculus of Constructions \cite{herbelin-1995-sequent-calculus}, but without inductive types or cyclic backlinks.
Our calculus extends Herbelin's work by adding:
\begin{itemize}
\item Inductive types and pattern matching
\item Cyclic backlinks (folding)
\item Focused phases (asynchronous vs synchronous)
\item Connection to supercompilation
\end{itemize}

\subsection{Induction: local rules vs global discharge}

\paragraph{Traditional induction principles.}
In standard type theory, induction is a derived principle from the elimination rule for inductive types.
To prove a property $P(x)$ for all $x : \mathsf{List}$, one applies the list eliminator (recursor) with a motive $P$ and proofs for the base and step cases.
This requires \emph{choosing the induction principle upfront}---a significant practical burden.

\paragraph{Cyclic induction as post-hoc discovery.}
Sprenger and Dam compare local well-founded induction rules to cyclic systems with global discharge conditions \cite{sprenger-dam-2003-circular-tree}.
They show these approaches are equivalent in expressive power but differ operationally: cyclic proofs discover the induction scheme during proof search rather than requiring it upfront.

Our work extends this insight to dependent types and connects it to supercompilation:
\begin{itemize}
\item The backlink corresponds to recognizing a recursive configuration
\item The trace condition corresponds to structural descent (termination)
\item The cycle structure \emph{is} the induction principle, read off from the proof graph
\end{itemize}

This flexibility is essential for advanced optimizations (fusion, deforestation, tupling) where the appropriate induction principle is not obvious from the problem statement.
